\documentclass[11pt,a4paper]{article}

% ===== XeLaTeX + Korean =====
\usepackage{kotex}
\usepackage{fontspec}
\setmainfont{Times New Roman}
\setmainhangulfont{AppleMyungjo}
\setsanshangulfont{Apple SD Gothic Neo}

% ===== Layout =====
\usepackage[a4paper,top=18mm,bottom=18mm,left=20mm,right=20mm]{geometry}
\usepackage{fancyhdr}
\setlength{\headheight}{24pt}
\usepackage{enumitem}
\usepackage{mdframed}
\usepackage{array}
\usepackage{booktabs}

\setlength{\parindent}{1em}
\linespread{1.18}

% ===== Header / Footer =====
\pagestyle{fancy}
\fancyhf{}
\renewcommand{\headrulewidth}{0.6pt}
\fancyhead[L]{\sffamily\bfseries 2026 고1 영어 변형 — 지문 38 정답 및 해설}
\fancyhead[R]{\sffamily 제스처와 기억}
\renewcommand{\footrulewidth}{0pt}
\fancyfoot[C]{\framebox[16mm][c]{\thepage}}

% ===== helpers =====
\newcommand{\qno}[1]{\par\vspace{10pt}\noindent{\sffamily\bfseries #1.}\ }
\newcommand{\instr}[1]{{\sffamily #1}\par\vspace{3pt}}

\newmdenv[linewidth=0.6pt,roundcorner=3pt,innertopmargin=7pt,innerbottommargin=7pt,
          innerleftmargin=9pt,innerrightmargin=9pt,skipabove=5pt,skipbelow=5pt]{pbox}

\begin{document}

\begin{center}
{\fontsize{15}{17}\selectfont\sffamily\bfseries [지문 38] 정답 및 해설}
\end{center}
\vspace{8pt}

% ===== 정답 표 =====
\begin{center}
\renewcommand{\arraystretch}{1.4}
\begin{tabular}{cccc}
\toprule
\textbf{문항} & \textbf{유형} & \textbf{정답} & \textbf{배점} \\
\midrule
1 & 문장 삽입 & ④ & 2점 \\
2 & 빈칸추론 & ② & 2점 \\
3 & 서술형 — 요지(한글) & (모범답안 참조) & 2점 \\
4 & 내용 불일치 & ④ & 2점 \\
\bottomrule
\end{tabular}
\end{center}

\vspace{10pt}

% ===== Q1 해설 =====
\qno{1}\textbf{[문장 삽입]} \textbf{정답: ④}

\par\vspace{4pt}
\textbf{삽입 문장:} ``This suggests that the act of gesturing itself, not just following instructions to gesture, plays a role in strengthening memory.''

\par\vspace{4pt}
\textbf{정답 근거:} ( ④ ) 앞 문장은 두 번째 실험(지시 없이 자발적으로 제스처를 한 경우)에서 ``We found the same patterns''라며 동일한 결과를 보고한다. 삽입 문장의 ``This suggests''는 바로 이 결과를 가리키며, 제스처 \textit{자체}(지시를 따른 것이 아닌)가 기억 강화에 역할을 한다는 결론을 이끌어 낸다. 이어지는 ( ⑤ ) 문장 ``Producing gesture along with speech makes the information encoded in that speech memorable.''은 이 결론을 일반화한 최종 주장이므로, 삽입 문장은 ④ 자리가 가장 자연스럽다.

\par\vspace{4pt}
\textbf{오답 소거:}
\begin{itemize}[leftmargin=2em,itemsep=2pt]
  \item ③ — 두 번째 실험을 \textit{소개}하는 문장 앞이므로, 아직 그 결과가 제시되지 않은 시점이다. ``This suggests''가 가리킬 선행 결과가 없어 부적절.
  \item ①②⑤ — 문장의 논리 흐름(첫 번째 실험 소개 → 결과 → 두 번째 실험 소개 → 결과 → 결론)에서 각각 너무 이른 위치이거나(①②), 이미 최종 결론이 내려진 뒤(⑤)여서 부적절.
\end{itemize}

% ===== Q2 해설 =====
\qno{2}\textbf{[빈칸추론]} \textbf{정답: ② of their own accord}

\par\vspace{4pt}
\textbf{정답 근거:} 빈칸이 포함된 문장은 두 번째 실험의 결과로, ``지시를 받지 않고(no instructions about gesture)'' 제스처를 한 경우에도 기억이 더 좋았다는 내용이다. 원문의 ``spontaneously(자발적으로)''를 대체하는 표현이므로 \textbf{of their own accord(자발적으로, 스스로)}가 정답이다.

\par\vspace{4pt}
\textbf{오답 소거:}
\begin{itemize}[leftmargin=2em,itemsep=2pt]
  \item ① in response to direct instruction — 지시에 반응하여: 본문이 명시적으로 이 상황의 \textit{반대}임을 설명(``no instructions about gesture''). 매력적 오답이지만 첫 번째 실험을 설명하는 표현.
  \item ③ at a rapid pace — 빠른 속도로: 속도와 관련 없음.
  \item ④ in a repetitive manner — 반복적으로: 반복 여부는 본문에서 다루지 않음.
  \item ⑤ with great exaggeration — 크게 과장하여: 과장 정도와 관련 없음.
\end{itemize}

% ===== Q3 해설 =====
\qno{3}\textbf{[서술형 — 요지(한글)]} \textbf{모범답안}

\par\vspace{6pt}
\begin{pbox}
제스처를 하며 말하면 그 내용이 기억에 더 잘 남는다.

\vspace{4pt}
\textit{또는:} 자발적인 제스처는 기억력 향상에 도움이 된다.
\end{pbox}

\par\vspace{6pt}
\textbf{채점 포인트} (각 1점, 총 2점):
\begin{itemize}[leftmargin=2em,itemsep=2pt]
  \item \textbf{[1점]} ``제스처'' 또는 이에 상응하는 표현이 포함될 것.
  \item \textbf{[1점]} ``기억 향상/기억에 잘 남음'' 또는 이에 상응하는 표현이 포함될 것.
\end{itemize}

\par\vspace{4pt}
\textbf{오답 처리:} 제스처와 기억의 관계를 언급하지 않거나, 40자를 초과하면 0점.

\vspace{8pt}
\qno{4}\textbf{[내용 불일치]} \textbf{정답: ④}
\par 두 번째 실험에서도 자발적으로 제스처를 한 항목을 그렇지 않은 항목보다 더 잘 기억했다는 동일한 패턴이 발견되었다. 따라서 기억 차이가 없었다는 ④는 본문과 일치하지 않는다.

\end{document}
